\section{Discussion and Limitations}
\subsection{What the Optimization Audit Establishes}
The experiment separates a content intervention from its ranking outcome. Applying a query-blind rewrite to seven fact-equivalent inputs does not make the resulting retrieval robust: the two adapted baselines increase VI, and their effectiveness varies by dataset and encoder. Canonicalization before rewriting blocks incoming-order variation under deterministic caching, but the subsequent rewrite can still reduce inclusion. Once a strict factuality guard is imposed, nearly every output falls back to canonical text. The defensible engineering result is consequently modest: normalize raw atoms before indexing, verify generated text against the complete fact inventory, and evaluate recall as well as stability.

This finding does not overturn E-GEO. Its published experiment rewrites one item within a fixed ten-candidate reranking set with GPT-4.1; ours adapts two released prompts to a 0.5B local model and asks whether products enter a dense first-stage pool. The poor fact retention and frequent 512-token truncation primarily characterize this adaptation. They nevertheless expose a deployment requirement that prompt-level instructions alone do not satisfy: a product rewriting pipeline needs explicit factuality and failure handling before rank changes can be interpreted as relevance-preserving visibility gains.

The generation control also changes what can be claimed causally. Raw permutations are verified by atom equality, so their target-only rank changes isolate representation. Free-text generations are not equivalent by construction. Under stochastic decoding, the same input sometimes crosses Top-100 across replicates as often as different input orders do. These sources must be modeled separately; pooling all outputs would over-attribute generator noise and semantic drift to serialization.

\subsection{From Retrieval Visibility to User Exposure}
Visibility here means eligibility at an offline retrieval cutoff. It is not observed exposure, clicks, satisfaction, purchase, or welfare, and relevance labels are not objective product quality. The shopping-agent motivation explains why candidate inclusion matters but does not substitute for an agent or user evaluation. The mechanism may arise wherever unordered records are serialized for retrieval, including jobs, hotels, restaurants, and events; neither its existence nor magnitude in those domains is established here.

\subsection{Limitations and Evaluation History}
The evidence covers two English public datasets, synthetic fact-preserving permutations, fixed catalogs, and exact offline retrieval. ESCI judgments are incomplete and its historical common catalog differs from official small candidate pools. The optimization audit contains only eight frozen queries per dataset; query-cluster intervals are accordingly wide, and the two-product common-valid subset supports no effectiveness inference. The source datasets and earlier test results had already been inspected before this post-hoc extension, although sample selection and methods were frozen before extension ranks were read.

The adapted generator differs from the original model, data, objective, and candidate setting. The automated fact checker is conservative, incomplete, and not human ground truth; source conflicts make WANDS especially demanding. Greedy decoding and deterministic caching are implementation choices. The stochastic control is small. The measured generation exceeded the planned eight-hour stop by about 3.5 hours across primary and noise runs; no method, sample, or stopping decision used ranking outcomes, but this remains a protocol deviation. We test no paid API, trained CHARM model, ANN index, commercial system, live shopping agent, or user behavior. Confidence intervals containing zero do not establish equivalence, and no fairness or disparate-impact conclusion follows from these results.
